% ============================================================================
% Chapter 9: File and Folder Management
% ============================================================================

\chapter{File and Folder Management}
\label{ch:files}

\begin{objectives}
  \item Understand the concept of files and folders
  \item Create, rename, copy, move, and delete files and folders
  \item Navigate the File Explorer window
  \item Know the rules for naming files in Windows
\end{objectives}

\bytetip[tip]{Think of your computer like a giant filing cabinet. \textbf{Folders} are the drawers, and \textbf{files} are the documents inside them. Keeping things organised will save you \emph{hours} of searching!}

% ---------------------------------------------------------------
\section{What Are Files and Folders?}
\label{sec:files-folders}
\index{file}\index{folder}

\begin{theorybox}[Files and Folders]
\begin{itemize}[leftmargin=*, label={\textcolor{ModuleBlue}{\faAngleRight}}, itemsep=4pt]
  \item A \textbf{file}\index{file} is a collection of data stored with a name and extension (e.g., \texttt{MyEssay.docx}, \texttt{Photo.jpg})
  \item A \textbf{folder}\index{folder} (also called a \textbf{directory}) is a container that holds files and other folders
  \item Folders help you \textbf{organise} your files so you can find them easily
\end{itemize}
\end{theorybox}

% --- Figure 9.1: File and Folder Hierarchy Tree ---
\begin{figure}[H]
\centering
\begin{tikzpicture}[
  level distance=1.4cm, sibling distance=3.5cm,
  folder/.style={draw=FunYellow!80!black, fill=FunYellow!20, rounded corners=3pt,
    font=\small\sffamily\bfseries, minimum height=0.7cm, minimum width=2cm, align=center},
  file/.style={draw=NavyBlue!50, fill=white, rounded corners=1pt,
    font=\small\sffamily, minimum height=0.6cm, minimum width=2cm, align=center},
  edge from parent/.style={draw=NavyBlue, line width=1.2pt, ->, >=Stealth}
]
  \node[folder] {Documents}
    child { node[folder] {School}
      child { node[file] {\faFileWord\ Maths.docx} }
      child { node[file] {\faFileWord\ Science.docx} }
    }
    child { node[folder] {Photos}
      child { node[file] {\faImage\ Holiday.jpg} }
    }
    child { node[file] {\faFile\ MyStory.txt} };
\end{tikzpicture}
\caption{Files and folders form a tree-like hierarchy}
\label{fig:file-tree}
\end{figure}

% ---------------------------------------------------------------
\section{File Operations}
\label{sec:file-ops}
\index{file!operations}

Here are the common things you can do with files:

% --- Figure 9.2: File Operations Step Map ---
\begin{figure}[H]
\centering
\begin{tikzpicture}[node distance=0.3cm]
  \node[process, fill=PartColor1!10, draw=PartColor1, minimum width=1.8cm] (s1) {Create};
  \node[process, fill=PartColor3!10, draw=PartColor3, minimum width=1.8cm, right=of s1] (s2) {Open};
  \node[process, fill=PartColor6!10, draw=PartColor6, minimum width=1.8cm, right=of s2] (s3) {Edit};
  \node[process, fill=PartColor4!10, draw=PartColor4, minimum width=1.8cm, right=of s3] (s4) {Save};
  \node[process, fill=PartColor2!10, draw=PartColor2, minimum width=2cm, right=of s4] (s5) {Copy/Move};
  \node[process, fill=PartColor5!10, draw=PartColor5, minimum width=1.8cm, right=of s5] (s6) {Delete};
  \foreach \a/\b in {s1/s2, s2/s3, s3/s4, s4/s5, s5/s6}{
    \draw[thickarrow, line width=1.2pt] (\a) -- (\b);
  }
\end{tikzpicture}
\caption{Common file operations in sequence}
\label{fig:file-ops}
\end{figure}

\begin{shortcutbox}
\begin{tabular}{ll}
  \key{Ctrl}+\key{C} & Copy selected file \\
  \key{Ctrl}+\key{X} & Cut (move) selected file \\
  \key{Ctrl}+\key{V} & Paste file \\
  \key{Delete} & Send file to Recycle Bin \\
  \key{F2} & Rename selected file \\
  \key{Ctrl}+\key{Z} & Undo last action \\
\end{tabular}
\end{shortcutbox}

% ---------------------------------------------------------------
\section{File Naming Rules}
\label{sec:naming-rules}
\index{file!naming rules}

% --- Table 9.1: Valid vs Invalid File Names ---
\begin{table}[H]
\centering
\caption{Valid vs invalid file names in Windows}
\label{tab:file-names}
\renewcommand{\arraystretch}{1.3}
\begin{tabular}{ll}
\toprule
\cellcolor{LabGreen!15}\textcolor{LabGreen}{\ding{51}} \textbf{Valid Names} &
\cellcolor{WarnRed!10}\textcolor{WarnRed}{\ding{55}} \textbf{Invalid Names} \\
\midrule
\rowcolor{RowLight}
MyReport & My/Report \footnotesize(slash not allowed) \\
\rowcolor{RowDark}
Science\_Notes\_2024 & File:Name \footnotesize(colon not allowed) \\
\rowcolor{RowLight}
Chapter 1 Draft & CON \footnotesize(reserved name in Windows) \\
\rowcolor{RowDark}
photo\_holiday & file* \footnotesize(asterisk not allowed) \\
\bottomrule
\end{tabular}
\end{table}

\bytetip[warn]{Never delete files from the \textbf{Windows} or \textbf{System32} folder. These contain critical files that your computer needs to work. Deleting them can \emph{break your entire computer}!}

\begin{notebox}
Windows file names \textbf{cannot contain} these characters: \texttt{\textbackslash\ / : * ? " < > |}. Also, names like CON, PRN, AUX, NUL are reserved by Windows and can't be used.
\end{notebox}

\begin{funfact}
The longest file name allowed in Windows is \textbf{255 characters} (including the extension). But the \emph{full path} (including all folder names) can't exceed about 260 characters. So if you nest folders very deeply with long names, you might run into problems!
\end{funfact}

\begin{reviewqs}
  \item What is the difference between a file and a folder?
  \item List three operations you can perform on a file.
  \item Name three characters that cannot be used in a Windows file name.
  \item What is the keyboard shortcut for copying a file?
  \item What happens when you press the \key{Delete} key on a selected file?
\end{reviewqs}

\begin{challenge}
Create the following folder structure on a computer: \texttt{MyProject > Documents}, \texttt{MyProject > Images}, \texttt{MyProject > Videos}. Inside Documents, create a text file called ``Notes.txt.'' Then copy it to the Images folder and rename the copy to ``ImageNotes.txt.''
\end{challenge}
